In this paper, we propose instead to use the viewpoint of Josuat-Verg\`es, Menous, Novelli, and Thibon put forward in \cite{JMNT}, where these authors introduce a Hopf algebra $\cHS(\A)$ of decorated Schr\"oder trees to rewrite formula \eqref{eq:freeCM} as a sum over of Schr\"oder trees. Our main result gives an analogous formula for expressing monotone cumulants in terms of the moments as a sum indexed by Schr\"oder trees. See Sections \ref{sec:schroder} and \ref{sec:monotone} for the precise definitions appearing in the formulas.

\begin{Theorem}\label{thm:mainthm1}
Let $(\A,\varphi)$ be a noncommutative probability space and let $\{h_n\colon \A^n\to\mathbb{C}\}_{n\geq1}$ be the family of monotone cumulants of $(\A,\varphi)$. Then, for any $n\geq1$ and elements $a_1,\dots,a_n\in \A$, we have that
\begin{equation}\label{eq:MonMom1}
	h_n(a_1,\dots,a_n) =\sum_{t\in \ST(n)} \omega(\sk(t))\varphi_{\pi(t)} (a_1,\dots,a_n),
\end{equation}
where $\ST(n)$ denotes the set of Schr\"oder trees with $n+1$ leaves and, for a Schr\"oder tree $t$, $\sk(t)$ denotes the so-called skeleton of $t$, and $\omega(\sk(t))$ is the Murua coefficient associated to~$\sk(t)$.
\end{Theorem}

We remark that \eqref{eq:MonMom1} may be rewritten in the form of \eqref{eq:objective} by regrouping the trees according to which $\pi\in \NC(n)$ satisfies that $\pi(t)=\pi$. Thus \eqref{eq:MonMom1} may be seen as a refinement of the desired formula~\eqref{eq:objective} (see Corollary~\ref{cor:maincor} for the precise statement).

Our approach is based on a series of recent papers by Ebrahimi-Fard and Patras \cite{EFP1,EFP2,EFP3,EFP4,EFP5}, where the authors have developed a Lie-theoretical framework for cumulants in noncommutative probability. In such a picture, the relations between moments and the different types of cumulants can be explained analogously to the relation between a Lie algebra and its corresponding (Lie) group. More precisely, given a noncommutative probability space $(\A,\varphi)$, we can consider a~Hopf algebra $H$ given by the double tensor algebra defined over $\A$, and a character $\Phi\colon H \to \mathbb{C}$ extending the linear functional $\varphi$. The free, Boolean, and monotone cumulant functionals can be identified with certain infinitesimal characters $\kappa$, $\beta$ and $\rho$ defined on $H$, respectively. The important feature of~$H$ is that its coproduct splits into two half-coproducts, providing two non-associative products $\prec$ and $\succ$ on the algebraic dual $H^*$, such that $(H^*,\prec,\succ)$ is a unital noncommutative shuffle algebra~\cite{EFP1}.

Within this framework, the relations between moments and free, Boolean, and monotone cumulants are encoded by the three exponential maps $\mathcal{E}_\prec$, $\mathcal{E}_\succ$ and $\exp^*$ associated to $\prec$, $\succ$ and the associative product $*$ of $H^*$, respectively:
\begin{gather*}
\Phi = \mathcal{E}_\prec(\kappa),
\qquad \Phi = \mathcal{E}_{\succ}(\beta),
\qquad \Phi = \exp^*(\rho).
\end{gather*}
We observe that in this context, the problem of expressing the monotone cumulants in terms of the moments is simply equivalent to computing the logarithm with respect to the associative product $\rho = \log^*(\Phi).$ However, finding a closed formula in terms of noncrossing partitions is not obvious from this relation.
